\subsubsection{An exact central-order obstruction for spatial comparison}
\label{sec:centralmomentratio}
The spatially dependent coefficient in
\eqref{eq:comparisonmissing} cannot be chosen
independently of Laguerre order. This is a
consequence of the positive theta convolution,
without dividing by any transform value away
from the centre. We record the exact
moment identity and the resulting obstruction.

\begin{theorem}[Uniform divergence of the central comparison ratio]
\label{thm:centralmomentratio}
Let $\Phi_t(u)=e^{tu^2}\Phi(u)$ for
$0\le t\le1/2$, and let
$H_t(z)=\int_0^\infty\Phi_t(u)\cos(zu)\,du$.
For $n\ge0$ write $L_{n,t}(x)$ for the
generalized Laguerre coefficient of
$H_t(x+iy)H_t(x-iy)$ and
$k_n=(n+1)(2n+1)$. Set
\[
g_t(w)=(\Phi_t*\Phi_t)(w)>0,\qquad
M_{2n}(t)=\int_{\mathbb R}w^{2n}g_t(w)\,dw.
\]
Then
\begin{equation}\label{eq:centralratio}
\boxed{\quad
\frac{k_nL_{n+1,t}(0)}{L_{n,t}(0)}
=\frac12\frac{M_{2n+2}(t)}{M_{2n}(t)}.
\quad}
\end{equation}
The right-hand side is strictly increasing
in $n$ and tends to infinity. In fact,
\begin{equation}\label{eq:centraluniform}
\boxed{\quad
\lim_{n\to\infty}\inf_{0\le t\le1/2}
\frac{k_nL_{n+1,t}(0)}{L_{n,t}(0)}
=+\infty.
\quad}
\end{equation}
Consequently no coefficient that remains
bounded uniformly in $n$ at $x=0$
can satisfy
$k_nU_{n+1}(s,x)\le c_n(s,x)U_n(s,x)$
for all orders and the entire heat-time
interval.
\end{theorem}

\begin{proof}
Since the Jacobi kernel is even, the
Fourier transform has the full-line form
\[
H_t(z)=\frac12\int_{\mathbb R}
\Phi_t(u)e^{izu}\,du.
\]
At $x=0$, the product
$H_t(iy)H_t(-iy)=H_t(iy)^2$
is the bilateral Laplace transform of
the positive convolution measure:
\[
H_t(iy)^2=\frac14\int_{\mathbb R}
g_t(w)e^{-yw}\,dw
=\sum_{n\ge0}
\frac{M_{2n}(t)}{4(2n)!}y^{2n}.
\]
Odd moments vanish by parity.
Thus $L_{n,t}(0)=M_{2n}(t)/(4(2n)!)>0$.
The factorial identity
\[
\frac{(n+1)(2n+1)(2n)!}{(2n+2)!}
=\frac12
\]
proves \eqref{eq:centralratio}.

Cauchy--Schwarz in the positive measure
$g_t(w)\,dw$ gives
\[
M_{2n+2}(t)^2\le
M_{2n}(t)M_{2n+4}(t).
\]
The inequality is strict because $g_t$
has a continuous strictly positive density
on the whole real line, rather than
support concentrated at a single value
of $w^2$. Dividing by the positive
moments shows strict increase in $n$.

To prove divergence, fix $A>0$ and choose
$B>A$. The positive measure has strictly
positive mass in $B\le|w|\le B+1$.
Under the probability distribution
with density proportional to
$w^{2n}g_t(w)$, the probability of
$|w|\le A$ is bounded by a constant
times $(A/B)^{2n}$ and tends to zero.
The corresponding expectation of
$w^2$ is
$M_{2n+2}(t)/M_{2n}(t)$, and its
limit inferior is therefore at least
$A^2$. Because $A$ is arbitrary,
the quotient diverges.

All $M_{2n}(t)$ are finite and
continuous in $t\in[0,1/2]$, by the
super-exponential decay of the theta
kernel and dominated convergence.
For every fixed threshold $C>0$,
the sets
\[
O_n=\left\{t\in[0,1/2]:
\frac{M_{2n+2}(t)}{2M_{2n}(t)}>C\right\}
\]
are open relative to the heat-time
interval and increasing in $n$.
Their union covers that compact interval
by pointwise divergence.
A finite subcover, followed by monotonicity,
gives an index $N$ such that
$O_N=[0,1/2]$. The same inequality holds
for all $n\ge N$, proving
\eqref{eq:centraluniform}.
\end{proof}

A proposed envelope of the form
$c_n(s,x)\le C\{1+\log^2(2+|x|)\}$,
with a constant $C$ independent of $n$,
is therefore \emph{impossible already at
$x=0$}. An order-dependent envelope
$c_n(s,x)\le C_n\{1+\log^2(2+|x|)\}$
is not excluded: it must simply have
$C_n\to\infty$. A quantitative estimate
on that required growth (for example,
one of order $\log^2(n+2)$ at the centre)
is a research candidate, not a theorem
proved here.

\paragraph{Exact and numerical checks.}
The finite algebra behind
\eqref{eq:centralratio} and the
Cauchy--Schwarz ratio comparison is
formalized in
\path{XiKernel/CentralMomentRatio.lean}.
The rational certificate
\path{certificates/xi_central_moment_ratio_exact.py}
checks the factorial cancellation and
moment log-convexity on an independent
positive discrete model.
For the genuine theta kernel, independent
Gauss--Legendre quadratures at 512, 1024
and 2048 nodes agree in the displayed
precision. The ratios for
$n=1,5,20,55$ are approximately
\[
\begin{array}{c|cccc}
n&1&5&20&55\\ \hline
t=0&.033451&.108858&.305956&.587252\\
t=1/2&.033781&.109741&.307551&.589104
\end{array}
\]
These computations illustrate growth
at finite order and are not directed
interval certificates of its asymptotic
rate. The monotonicity and divergence
in Theorem~\ref{thm:centralmomentratio}
follow analytically, not from
these sampled values.

The central identity also pinpoints a
limitation of this route to RH:
$L_{n,t}(0)>0$ holds for the full theta
heat family regardless of whether
$H_0$ has any nonreal zeros. The missing
global proof must therefore use
noncentral information, especially
where Fourier oscillations can reverse
the sign of a Laguerre expression.
